Annual 10-K filings help separate more credible AI disclosure from narrative positioning. The classification of AI-related disclosure in 10-K filings is scalable and based on audited reporting content. The main construct, PatentMismatch, combines low-credibility disclosure composition with contemporaneous patent weakness. In the 2016--2025 sample, low-credibility AI disclosure rises sharply after 2022; for a substantial subset of firms, AI language increases more rapidly than observable patenting activity.

Across alternative definitions, the filing-based mismatch measure predicts weaker subsequent AI patent grants. This specification is the most robust version of the construct across tests. In the AI-talking sample, mismatch firm-years subsequently exhibit lower profitability and higher R\&D intensity. The patterns are more consistent with delayed realization and subsequent adjustment in real activity than with purely non-informative disclosure.

Low-credibility disclosure also varies with scrutiny and incentives. Among continuing AI disclosers, disclosure composition becomes more credible following SEC comment-letter scrutiny, while AI Focus remains largely unchanged. The SEC's March 18, 2024 AI-washing enforcement actions are followed by similar disclosure-cleanup patterns among pre-exposed firms (SEC, 2024). The March 2024 evidence is better read as regulatory salience or acceleration than as a sharply identified enforcement effect. Low-credibility AI disclosure is also more common when CEO equity incentives are stronger. Around large equity-issuance windows, disclosure credibility deteriorates prior to issuance and partly reverts afterward. Low-credibility AI disclosure appears shaped by both monitoring and opportunistic incentives.

The market-pricing results are more limited. Filing-date returns do not cleanly separate mismatch from non-mismatch firms, and post-filing return evidence remains benchmark-sensitive. The strongest return result is a short-horizon equal-weight alpha, with no comparable value-weight evidence. The issue-window evidence suggests a state-dependent rather than uniform short-run return pattern. Market tests therefore support a narrower interpretation than a broad AI-washing anomaly. The stronger evidence is that the filing-based credibility measure predicts later technological realization and organizational behavior, even when market pricing remains incomplete and uneven.

AI-related narratives can diffuse faster than observable implementation. That gap matters when investors and regulators must evaluate claims about technologies whose adoption is costly, valuable, and difficult to verify. The same disclosure-credibility problem can arise in other high-salience domains where claims become cheap before implementation is observable.
